\documentclass[12pt,a4paper]{article}
\usepackage[a4paper,margin=2.2cm]{geometry}
\usepackage{fontspec}
\usepackage{polyglossia}
\setmainlanguage{english}
\setotherlanguage{russian}
\setmainfont{Noto Serif}
\setsansfont{Noto Sans}
\setmonofont{DejaVu Sans Mono}
\usepackage{amsmath,amssymb,mathtools}
\usepackage{array,longtable,booktabs}
\usepackage{enumitem}
\usepackage{hyperref}
\hypersetup{colorlinks=true,linkcolor=black,urlcolor=blue,citecolor=black}
\setlist{nosep,leftmargin=*}
\DeclareMathOperator{\Spin}{Spin}
\DeclareMathOperator{\Tr}{Tr}
\newcommand{\TMI}{\mathcal T_{MI}}
\newcommand{\Adm}{\operatorname{Adm}}
\newcommand{\Struct}{\operatorname{Struct}}
\newcommand{\Dist}{\operatorname{Dist}}
\newcommand{\Cond}{\operatorname{Cond}}
\newcommand{\SM}{\mathrm{SM}}
\newcommand{\GR}{\mathrm{GR}}
\newcommand{\QM}{\mathrm{QM}}
\newcommand{\QG}{\mathrm{QG}}
\newcommand{\BH}{\mathrm{BH}}
\newcommand{\DE}{\mathrm{DE}}
\newcommand{\DM}{\mathrm{DM}}
\newcommand{\dd}{\,d}
\title{TMI-ToE 1.0\\\large Interface-Field Architecture for the Unification of Fundamental Interactions}
\author{Theory of Manifold Interfaces}
\date{Version 1.0}
\begin{document}
\maketitle
\tableofcontents
\newpage

\section*{Abstract}
\addcontentsline{toc}{section}{Abstract}

This document fixes the module \textbf{TMI-ToE 1.0} within the Theory of Manifold Interfaces. Its purpose is not to claim a completed ``Theory of Everything'', but to define a minimal consistent architecture in which the four fundamental interactions are described as projections of a unified interface curvature.

Short status:
\[
\boxed{\mathrm{TMI\mbox{-}ToE\ 1.0}\neq \text{a complete Theory of Everything}}
\]
\begin{center}
\fbox{\parbox{0.88\linewidth}{\centering $\mathrm{TMI\mbox{-}ToE\ 1.0}$ is a minimal interface-field architecture of unification.}}
\end{center}

Main idea:
\[
\boxed{F_a=\pi_a(\mathbb F_I),\qquad a\in\{G,S,W,EM\}}
\]
where $\mathbb F_I$ is the universal interface curvature, and $F_G,F_S,F_W,F_{EM}$ are the gravitational, strong, weak, and electromagnetic regimes.

\section{Initial status of the module}

The physical branch of TMI already contains several modules:
\begin{itemize}
  \item \textbf{TMI-QG}: interface geometric decoherence, the quantity $\Gamma_I$, a testable hypothesis, and a statistical form.
  \item \textbf{TMI-Dark}: the dark sector as a regime of the interface field $\chi_I=\chi_0+\delta\chi$.
  \item \textbf{TMI-BH}: a black hole as a one-way causal-informational interface.
  \item \textbf{TMI-ToE}: an architecture unifying these physical layers through a universal connection, curvature, and interface field.
\end{itemize}

TMI-ToE does not replace established theories. It requires general relativity, the Standard Model, and quantum mechanics to be recovered as limiting cases.

\section{TMI-ToE-0.9: consistency conditions}

The goal of version 0.9 is not to add new concepts, but to close the technical conditions of admissibility. If any of these conditions fails, the current TMI-ToE version must be revised.

\subsection{Recovery of known theories}

\subsubsection*{General-relativistic limit}

If the interface sector is switched off:
\[
\chi_I\to0,\qquad \mathcal L_{mix}\to0,
\]
then:
\[
\boxed{S_{TMI-ToE}\to S_{GR}+S_{SM}}.
\]
If matter and gauge fields are also removed:
\[
\mathcal L_{SM}\to0,
\]
then:
\[
\boxed{S_{TMI-ToE}\to S_{GR}}.
\]

\subsubsection*{Standard Model limit}

For a fixed geometry:
\[
g_{\mu\nu}=g_{\mu\nu}^{fixed}
\]
and vanishing interface mixing:
\[
\chi_I\to0,\qquad \zeta_a,\eta_H,y_I\to0,
\]
one must obtain:
\[
\boxed{S_{TMI-ToE}\to S_{SM}}.
\]
Thus the low-energy gauge limit must contain
\[
\boxed{SU(3)\times SU(2)\times U(1)}.
\]

\subsubsection*{Quantum limit}

For a fixed background and a weak interface sector, standard quantum dynamics must be recovered:
\[
\boxed{i\hbar\frac{\partial}{\partial t}\Psi=\widehat H\Psi}
\]
or the corresponding quantum field-theoretic form on a fixed background.

\subsection{The universal group $G_I$}

The minimal requirement is:
\[
\boxed{G_I\supset \Spin(1,3)\times SU(3)\times SU(2)\times U(1)}.
\]

For version 1.0, the safest form is not to assert a new simple group, but to fix an inclusion structure:
\[
\boxed{G_I=G_{grav}\ltimes G_{SM}\ltimes G_{\chi}}
\]
where
\[
G_{grav}\supset\Spin(1,3),\qquad G_{SM}=SU(3)\times SU(2)\times U(1).
\]

Thus $G_I$ is the group of admissible interface transformations containing gravitational, gauge, and interface symmetries.

\subsection{Projection consistency}

We introduce algebra projections:
\[
\boxed{\pi_a:\mathfrak g_I\to\mathfrak g_a,\qquad a\in\{G,S,W,EM\}},
\]
where $\mathfrak g_I=\mathrm{Lie}(G_I)$.

Then the physical curvatures are:
\[
\boxed{F_G=\pi_G(\mathbb F_I)},\qquad
\boxed{F_S=\pi_S(\mathbb F_I)},
\]
\[
\boxed{F_W=\pi_W(\mathbb F_I)},\qquad
\boxed{F_{EM}=\pi_{EM}(\mathbb F_I)}.
\]

\subsection{Gauge invariance}

The mixing term must preserve the gauge structure of the Standard Model. A safe form is:
\[
\boxed{
\mathcal L_{mix}^{gauge}
=-\frac14\sum_{a=1}^{3}\zeta_a\frac{\chi_I^2}{M_I^2}\Tr(F_{a\mu\nu}F_a^{\mu\nu})
}
\]
which is admissible if $\chi_I$ is a gauge singlet:
\[
\boxed{\chi_I\sim(1,1,0)}.
\]

\subsection{Fermionic consistency}

A raw term $\chi_I\bar\Psi\Psi$ may violate gauge invariance before electroweak symmetry breaking. A safer coupling to fermions is:
\[
\boxed{
\mathcal L_{\chi\Psi}
=
\sum_f y_{I,f}\frac{\chi_I}{M_I}\bar\Psi_{L,f}H\Psi_{R,f}+h.c.
}
\]
where $H$ is the Higgs field and $f$ indexes fermion type and generation.

\subsection{Higgs sector}

A compatible potential is:
\[
\boxed{
V(H,\chi_I)=
-\mu_H^2H^\dagger H+
\lambda_H(H^\dagger H)^2+
V(\chi_I)+
\eta_H\chi_I^2H^\dagger H
}
\]
where
\[
\boxed{V(\chi_I)=V_0+\frac12m_I^2\chi_I^2+\frac{\beta_I}{4}\chi_I^4}.
\]

The effective Higgs parameter is:
\[
\boxed{\mu_{H,eff}^2=\mu_H^2-\eta_H\chi_I^2}.
\]
Thus the Higgs vacuum expectation value may receive an interface correction:
\[
\boxed{v_H^{eff}=v_H(\chi_I)}.
\]

\subsection{Anomaly cancellation}

If $\chi_I$ is a gauge singlet and no new chiral fermions are added, the anomaly structure of the Standard Model remains unchanged:
\[
\boxed{\mathcal A_{TMI-ToE}=\mathcal A_{SM}=0}.
\]
If new fermions or new $G_I$ representations are introduced in future versions, one must require:
\[
\boxed{\mathcal A_{G_I}=0}.
\]

\subsection{EFT status}

The present version must be treated as an effective interface-field theory:
\[
\boxed{TMI\mbox{-}ToE=\text{effective interface-field theory}}
\]
valid at energies
\[
\boxed{E\ll M_I}.
\]
Operators such as
\[
\frac{\chi_I^2}{M_I^2}F_{\mu\nu}F^{\mu\nu},\qquad
\frac{\chi_I}{M_I}\bar\Psi_LH\Psi_R
\]
are effective operators suppressed by the interface scale $M_I$.

\subsection{Stability}

The potential must be bounded from below. Sufficient conditions are:
\[
\boxed{\beta_I>0,\qquad \lambda_H>0}.
\]
The mixed potential must not generate an unstable vacuum. A working condition is:
\[
\boxed{\eta_H>-2\sqrt{\lambda_H\beta_I/4}}.
\]

\subsection{Compatibility with other TMI modules}

\subsubsection*{Connection with TMI-QG}
\[
\boxed{\chi_I\Rightarrow\Gamma_I},\qquad
\boxed{\Delta\Gamma_I=\kappa_Ix_I}.
\]

\subsubsection*{Connection with TMI-Dark}
\[
\boxed{\chi_I=\chi_0+\delta\chi},\qquad
\boxed{\chi_0\Rightarrow DE,\quad \delta\chi\Rightarrow DM}.
\]

\subsubsection*{Connection with TMI-BH}
On a black-hole horizon, the restrictions
\[
\boxed{\mathbb A_I|_{\mathcal H_g}},\qquad
\boxed{\mathbb F_I|_{\mathcal H_g}}
\]
must have a meaningful interpretation.

\section{TMI-ToE-1.0: passport of the minimal completed architecture}

\subsection{Name}

\[
\boxed{TMI\mbox{-}ToE\ 1.0}
\]

Full title:
\begin{center}
\fbox{\parbox{0.88\linewidth}{\centering Interface-Field Architecture for the Unification of Fundamental Interactions}}
\end{center}

\subsection{Status}

\begin{center}
\fbox{\parbox{0.88\linewidth}{\centering A minimal completed architecture, but not a complete physical Theory of Everything.}}
\end{center}

That is:
\[
\boxed{TMI\mbox{-}ToE\ 1.0\neq\text{a complete Theory of Everything}}.
\]

\subsection{Main object}

\[
\boxed{
\mathcal M_I^{ToE}=
(\mathcal M_U,\mathcal E_I,G_I,\mathbb A_I,\mathbb F_I,\Psi,H,\chi_I,\Adm_I,\Struct_I)
}
\]

\begin{longtable}{>{\raggedright\arraybackslash}p{0.23\linewidth}>{\raggedright\arraybackslash}p{0.68\linewidth}}
\toprule
Symbol & Meaning \\
\midrule
$\mathcal M_U$ & spacetime \\
$\mathcal E_I$ & universal interface bundle \\
$G_I$ & universal interface group \\
$\mathbb A_I$ & universal interface connection \\
$\mathbb F_I$ & universal interface curvature \\
$\Psi$ & full quantum-field sector \\
$H$ & Higgs field \\
$\chi_I$ & interface field \\
$\Adm_I$ & set of admissible physical transitions \\
$\Struct_I$ & structured physical result \\
\bottomrule
\end{longtable}

\subsection{Universal bundle}

\[
\boxed{\pi_I:\mathcal E_I\to\mathcal M_U}
\]

The fiber over a point is:
\[
\boxed{\mathcal F_{I,x}=\pi_I^{-1}(x)}.
\]

Working decomposition:
\[
\boxed{
\mathcal F_{I,x}=\mathcal F_G\oplus\mathcal F_{SU(3)}\oplus\mathcal F_{SU(2)}\oplus\mathcal F_{U(1)}\oplus\mathcal F_\Psi\oplus\mathcal F_H\oplus\mathcal F_\chi
}
\]

\subsection{Universal connection}

\[
\boxed{\mathbb A_I\in\Omega^1(\mathcal M_U,\mathfrak g_I)}.
\]

Working form:
\[
\boxed{
\mathbb A_I=
\omega^{ab}J_{ab}+e^aP_a+G^AT_A+W^iT_i+BY+A_\chi\Xi
}
\]
where $A_\chi$ is an interface one-form. If $\chi_I$ is taken as a scalar field, one may set:
\[
\boxed{A_\chi=d\chi_I}.
\]

\subsection{Universal curvature}

\[
\boxed{\mathbb F_I=d\mathbb A_I+\mathbb A_I\wedge\mathbb A_I}.
\]

Projections:
\[
\boxed{F_a=\pi_a(\mathbb F_I),\qquad a\in\{G,S,W,EM\}}.
\]

Main formula:
\begin{center}
\fbox{\parbox{0.88\linewidth}{\centering The four fundamental interactions are projections of a unified interface curvature.}}
\end{center}

\subsection{Unified action}

\[
\boxed{S_{TMI-ToE}=S_{GR}+S_{SM}+S_\chi+S_{mix}}.
\]

Expanded form:
\[
\boxed{
S_{TMI-ToE}=\int_{\mathcal M_U}\sqrt{-g}\left[
\frac{c^4}{16\pi G}(R-2\Lambda)+\mathcal L_{SM}+\mathcal L_\chi+\mathcal L_{mix}
\right]\dd^4x
}
\]

Interface Lagrangian:
\[
\boxed{
\mathcal L_\chi=
-\frac12g^{\mu\nu}\nabla_\mu\chi_I\nabla_\nu\chi_I
-V(\chi_I)-\frac12\xi_IR\chi_I^2
}
\]

Mixing term:
\[
\boxed{
\mathcal L_{mix}=
-\frac14\sum_{a=1}^{3}\zeta_a\frac{\chi_I^2}{M_I^2}\Tr(F_{a\mu\nu}F_a^{\mu\nu})
+\eta_H\chi_I^2H^\dagger H
+\sum_f y_{I,f}\frac{\chi_I}{M_I}\bar\Psi_{L,f}H\Psi_{R,f}+h.c.
}
\]

\subsection{Main chain}

\[
\boxed{
\mathcal E_I\Rightarrow G_I\Rightarrow\mathbb A_I\Rightarrow\mathbb F_I\Rightarrow\{F_G,F_S,F_W,F_{EM}\}\Rightarrow\Adm_I\Rightarrow\mathcal H_I\Rightarrow\Psi\Rightarrow\Dist_{poss}\Rightarrow\Struct_I
}
\]

Meaning:
\begin{enumerate}
  \item the unified interface structure defines admissible interactions;
  \item interactions define physical transitions;
  \item the quantum state defines the distribution of possibilities;
  \item the interface transforms possibilities into structure.
\end{enumerate}

\section{Testable consequences}

\subsection{Consequence 1: corrections to gauge couplings}

From the gauge mixing term
\[
-\frac14\zeta_a\frac{\chi_I^2}{M_I^2}\Tr(F_{a\mu\nu}F_a^{\mu\nu})
\]
one obtains the effective correction:
\[
\boxed{
\frac{1}{g_{a,obs}^2(E)}=
\frac{1}{g_{a,SM}^2(E)}+
\zeta_a\frac{\chi_I^2(E)}{M_I^2}
}
\]
and therefore:
\[
\boxed{
\Delta_a(E)=
\frac{1}{g_{a,obs}^2(E)}-
\frac{1}{g_{a,SM}^2(E)}=
\zeta_a\frac{\chi_I^2(E)}{M_I^2}
}
\]

Null hypothesis:
\[
\boxed{H_0:\zeta_a=0}.
\]
TMI-ToE hypothesis:
\[
\boxed{H_{TMI-ToE}:\exists a,\ \zeta_a\neq0}.
\]

\subsection{Consequence 2: corrections to the Higgs sector}

Through $V(H,\chi_I)$ one obtains:
\[
\boxed{\mu_{H,eff}^2=\mu_H^2-\eta_H\chi_I^2}
\]
and the possible test line:
\[
\boxed{\Delta v_H=v_H^{obs}-v_H^{SM}=F(\eta_H,\chi_I)}.
\]

\subsection{Consequence 3: relation to interface decoherence}

If:
\[
\Gamma_I=\lambda_I\|\chi_I\|^2
\]
and:
\[
\Delta_a(E)=\zeta_a\frac{\chi_I^2(E)}{M_I^2},
\]
then:
\[
\boxed{\Delta_a(E)=\frac{\zeta_a}{\lambda_I M_I^2}\Gamma_I(E)}.
\]

This is a cross-module prediction: the same interface intensity may appear both in decoherence and in gauge-coupling corrections.

\section{Criteria of support and weakening}

\subsection{Support criteria}

The module receives support if:
\begin{enumerate}
  \item $GR$, $SM$, and $QM$ are correctly recovered;
  \item the mixing terms preserve gauge symmetries;
  \item no uncancelled anomalies appear;
  \item the effects are suppressed by $M_I$ and do not contradict data;
  \item observed corrections share a common dependence on $\chi_I^2$.
\end{enumerate}

\subsection{Weakening criteria}

TMI-ToE weakens if:
\begin{enumerate}
  \item $GR$ or $SM$ is not recovered;
  \item $\mathcal L_{mix}$ violates gauge invariance;
  \item uncancelled anomalies appear;
  \item effective corrections contradict experimental constraints;
  \item no distinguishable testable consequence exists.
\end{enumerate}

\section{Final fixation}

\begin{center}
\fbox{\parbox{0.88\linewidth}{\centering $TMI\mbox{-}ToE\ 1.0$ is closed as a minimal consistent architecture.}}
\end{center}

But:
\begin{center}
\fbox{\parbox{0.88\linewidth}{\centering TMI-ToE 1.0 is not a completed physical Theory of Everything.}}
\end{center}

Proper formulation:
\begin{quote}
TMI-ToE-1.0 is a completed minimal architectural module of unification in which the four interactions are represented as projections of a unified interface curvature, while the interface field $\chi_I$ connects geometry, gauge fields, quantum state, dark sector, black holes, and the process of structuring.
\end{quote}

Shortest final formula:
\[
\boxed{\mathbb F_I=d\mathbb A_I+\mathbb A_I\wedge\mathbb A_I}
\]
\[
\boxed{F_a=\pi_a(\mathbb F_I),\qquad a\in\{G,S,W,EM\}}
\]
\[
\boxed{S_{TMI-ToE}=S_{GR}+S_{SM}+S_\chi+S_{mix}}
\]
\[
\boxed{\chi_I\Rightarrow\{\Gamma_I,\ DarkSector,\ \Delta_a(E),\ BH\mbox{-}interface\}}
\]

\end{document}
